\label{mab:rm:ch:chaos}

\section{Intrinsic construction of the dodecaphasic family}

The Feigenbaum constants characterize the universality of period
doubling. Their derivation in HMT first requires constructing
an analytic unimodal family, proving its quadratic criticality, and
then studying the action of the renormalization operator.

The initial map does not follow from an external harmonic choice. The return
of \(\TPK\) determines the oriented dodecaphase wheel
\(U_{12}^{\rm sign}\). In the uniform sector \(\eta=0\), the normalized
trace of its twelve sectors fixes \(w_m=1/12\), and the phase functional
\begin{equation}
\mathcal P_{\rm crit}:U_{12}^{\rm sign}\longrightarrow\R/2\pi\Z,
\qquad
m\longmapsto \frac{(30m-10)\pi}{180}
=\pi\frac{3m-1}{18}
\label{mab:rm:eq:critical-phase-reader}
\end{equation}
preserves order and orientation. The causal composition is
\[
\APP\to\TRIT\to\TPK\to U_{12}^{\rm sign}
\xrightarrow{\mathcal P_{\rm crit}}
\{(\phi_m,w_m)\}_{m=1}^{12}
\to u_0\to f_\lambda.
\]
Every modification of \(\mathcal P_{\rm crit}\), of the weights, or of the order of
the phases alters the family and constitutes a criterion of refutation. These
choices must be fixed before the contrast with \(\delta_F\).

The resulting dodecaphasic family starts from

\[
a_m=\frac{3m-1}{18},\qquad
\phi_m=\pi a_m,\qquad m=1,\ldots,12,
\]

and of the second moment

\[
D_0=\frac1{12}\sum_{m=1}^{12}a_m^2.
\]

The identity

\begin{equation}
\sum_{m=1}^{12}(3m-1)^2=5394=6\cdot899
\label{mab:rm:eq:chaos-sum2}
\end{equation}

yields

\[
D_0=\frac{899}{648},\qquad
s_0=\frac{36}{\pi\sqrt{899}}.
\]

Therefore,

\begin{equation}
s_0\phi_m=\frac{2(3m-1)}{\sqrt{899}}.
\label{mab:rm:eq:pi-cancellation}
\end{equation}

The factor \(\pi\) cancels before the iterations are performed. The family does not
depend on \(\delta_F\) or \(\alpha_F\) as input data.

Cancellation is a structural property of
normalization. Before normalization, the second angular moment
is
\[
\frac1{12}\sum_{m=1}^{12}\phi_m^2
=\pi^2D_0.
\]
The condition of quadratic criticality fixes
\(s_0^2\pi^2D_0=2\), and therefore each normalized phase depends only on the
integer arithmetic of \(3m-1\). The functional preserves the angular provenance of
the wheel, but the iteration receives an intrinsic real family. This is
precisely what makes it possible to distinguish a derivation from a calibration
by \(\pi\).

\section{Exact analytic representation and closed trigonometric expression}

We define

\begin{equation}
u_0(x)=1-\frac1{12}\sum_{m=1}^{12}
\cos\!\left(\frac{2(3m-1)}{\sqrt{899}}x\right).
\label{mab:rm:eq:chaos-profile}
\end{equation}

The development at the origin is

\begin{equation}
u_0(x)=x^2-\frac{715769}{2424603}x^4+O(x^6).
\label{mab:rm:eq:chaos-taylor}
\end{equation}

The finite sum admits the following closed expression. If \(y=x/\sqrt{899}\),

\begin{equation}
u_0(x)=1-\frac{\cos(37y)\sin(36y)}{12\sin(3y)}
=1-\frac{\sin(73y)-\sin y}{24\sin(3y)}.
\label{mab:rm:eq:chaos-closed}
\end{equation}

The apparent singularity at \(y=0\) is removable and reproduces \cref{mab:rm:eq:chaos-taylor}.

\begin{proof}[Derivation of the closed form]
The twelve arguments form an arithmetic progression:
\[
2(3m-1)y=(6m-2)y,\qquad m=1,\ldots,12.
\]
The sum of cosenos of a progresion satisfies
\[
\sum_{m=1}^{N}\cos(a+md)
=\frac{\sin(Nd/2)}{\sin(d/2)}
 \cos\!\left(a+\frac{N+1}{2}d\right).
\]
Taking \(N=12\), \(d=6y\) and \(a=-2y\) yields
\[
\sum_{m=1}^{12}\cos((6m-2)y)
=\frac{\sin36y}{\sin3y}\cos37y.
\]
The product--to--sum identity
\(2\cos37y\sin36y=\sin73y-\sin y\) gives the second expression for
\cref{mab:rm:eq:chaos-closed}.
\end{proof}

The quartic coefficient is not fitted either. The integer sums
\begin{equation}
\sum_{m=1}^{12}(3m-1)^2=5394,\qquad
\sum_{m=1}^{12}(3m-1)^4=4294614
\label{mab:rm:eq:chaos-moments24}
\end{equation}
inserted in the development of \(\cos z\) produce exactly
\[
\frac1{24\cdot12}\sum_m
\left(\frac{2(3m-1)}{\sqrt{899}}\right)^4
=\frac{715769}{2424603}.
\]
Thus, both the quadratic normalization and the first nonlinear
correction are finite moments of the same wheel.

\section{Pertenencia to the clase analitica unimodal cuadratica}

Consider we

\[
f_\lambda(x)=1-\lambda u_0(x),\qquad \lambda>0.
\]

For \(0<x\le1\), all arguments

\[
\frac{2(3m-1)}{\sqrt{899}}x
\]

belong to \((0,\pi)\). By therefore,

\[
u_0'(x)>0,\qquad u_0'''(x)<0.
\]

Since \(u_0\) is even, \(x=0\) is its unique critical point and \(u_0''(0)=2\). The Schwarzian of \(f_\lambda\) is

\begin{equation}
S(f_\lambda)
=\frac{u_0'''}{u_0'}
-\frac32\left(\frac{u_0''}{u_0'}\right)^2<0
\qquad(0<|x|\le1).
\label{mab:rm:eq:negative-schwarzian}
\end{equation}

\begin{theorem}[Exact membership in the unimodal class]
For
\[
0<\lambda\le \frac{2}{u_0(1)},
\]
the family \(f_\lambda=1-\lambda u_0\) is a self-map that is real--analytic on \([-1,1]\) and even, with a single critical point of quadratic type at the origin and negative Schwarzian away from that point. Consequently, without using Feigenbaum values as a target value, it belongs to the \(S\)-unimodal quadratic class subject to the action of the period-doubling renormalizer \cite{rm:Feigenbaum1978,rm:Lanford1982}. For \(\lambda\) outside that interval, the analytic germ and its local form are retained, but one does not automatically assert a self-map of the interval.
\end{theorem}

The theorem establishes the contribution intrinsic to HMT prior to the
universality theorem: the analytic germ and its type of criticality.

\begin{proof}[Detalles of the signo and of the self-map]
Deriving \cref{mab:rm:eq:chaos-profile},
\begin{align*}
u_0'(x)&=\frac1{12}\sum_m k_m\sin(k_mx),\\
u_0'''(x)&=-\frac1{12}\sum_m k_m^3\sin(k_mx),
\qquad k_m=\frac{2(3m-1)}{\sqrt{899}}.
\end{align*}
For \(0<x\le1\), \(0<k_mx<\pi\), so all summands in the first
line are positive and those in the second are negative. Parity extends
monotonicity to both sides of the origin. Since \(u_0(0)=0\) and
\(0\le u_0(x)\le u_0(1)\), the bound
\(0<\lambda\le2/u_0(1)\) implies
\(-1\le1-\lambda u_0(x)\le1\). Finally,
\(u_0''(0)=1/12\sum k_m^2=2\), so the critical point is
quadratic.
\end{proof}

\section{Sequences finite of parameters superstable}

The sucesion superestable certified until level eight produces

\[
\delta_8=4.66921218915\ldots,
\qquad
\alpha_8=2.50296847988\ldots.
\]

Its status is \emph{finite certificate free of target values}: each
approximant is obtained from the previously defined family, without introducing
the universal limit as a datum. Decimal agreement provides a
convergence indicator but does not replace the hyperbolicity theorem.

An independent numerical extension of the same sequence, with a step size
adapted to the separation between roots and control of the exact period,
provisionally reaches:

\begin{center}
\small
\begin{tabular}{@{}ccll@{}}
\toprule
Level & \(\lambda_n\) & \(\delta_n\) & \(\alpha_n\)\\
\midrule
11 & 1.87403828195439908045 & 4.66920170694424982745 & 2.50290847517165181764\\
12 & 1.87403836411023460451 & 4.66920163036308158848 & 2.50290800379001546915\\
13 & 1.87403838170549870372 & 4.66920161361839219913 & 2.50290790268748193108\\
14 & 1.87403838547386545431 & 4.66920161007472591023 & 2.50290788100969754274\\
\bottomrule
\end{tabular}
\end{center}

These levels have the status of \emph{New numerical certificate
pending reproducible registration}; they do not supersede the persistent certificate
through level eight. An Aitken extrapolation yields
\(\lambda_\infty^{\HMT}\simeq1.8740383865008918080\) and
\(\alpha_F\simeq2.50290787509\) as candidate limits, conditional on
a proof of hyperbolicity and transversality.

To avoid an unspecified numerical procedure, the method is formulated as
follows. Define
\begin{equation}
F_n(\lambda)=f_\lambda^{\,2^n}(0).
\label{mab:rm:eq:superstable-equation}
\end{equation}
The raiz superestable \(\lambda_n\) is the primera raiz of \(F_n\) to the
right of \(\lambda_{n-1}\) that satisfies simultaneously
\[
f_{\lambda_n}^{\,2^j}(0)\ne0
\quad(0\le j<n).
\]
This inequality excludes divisor periods. With disjoint intervals
\(I_n\ni\lambda_n\), one computes
\begin{equation}
\delta_n=
\frac{\lambda_{n-1}-\lambda_{n-2}}
     {\lambda_n-\lambda_{n-1}},
\qquad
\alpha_n=-\frac{d_{n-1}}{d_n},
\label{mab:rm:eq:feigenbaum-approximants}
\end{equation}
where \(d_n=f_{\lambda_n}^{\,2^{n-1}}(0)\) is the oriented distance from the
last new point to the critical point. A persistent certificate must retain
the root intervals, the exclusion of smaller periods, and the propagation
of error from \cref{mab:rm:eq:feigenbaum-approximants}; a mere printed decimal does not
satisfy those three conditions.

The discretization rule also constitutes a negative control. A
fixed lower bound greater than the separation
\(|\lambda_n-\lambda_{n-1}|\) can omit the first root and nevertheless produce a numerically plausible
succession. For this reason the increment
must scale with the previous separation and each candidate must be verified
through its exact period.

\section{Conditions for certifying the renormalization operator}

On a space of even analytic functions analytic, we define

\begin{equation}
\mathcal R(f)(x)=-a(f)^{-1}f\bigl(f(-a(f)x)\bigr),
\qquad a(f)=-f(1).
\label{mab:rm:eq:renormalization}
\end{equation}

The complete closure of \(\delta_F\) and \(\alpha_F\) requires jointly certifying:

\begin{enumerate}
\item existence of a fixed point \(g=\mathcal R(g)\) in a declared analytic ball;
\item uniqueness local of ese point fixed;
\item a single unstable real eigenvalue of \(D\mathcal R_g\), equal to \(\delta_F\);
\item spectral radius strictly less than one on the stable complement;
\item transversality of the curve \(\lambda\mapsto f_\lambda\) with respect to the stable manifold of \(g\).
\end{enumerate}

A natural HMT implementation employs Chebyshev truncations of order
\(9^m\), interval arithmetic, and radii polynomials. The proof must
provide a ball, an approximate inverse, a defect bound, and a
Lipschitz bound satisfying \(Y+Z(r)<r\). Thus, the nonadic
depth parameterizes a fixed-point certificate and not a mere sequence
of decimal approximations.

On a Chebyshev even basis,
\[
g(x)=\sum_{j=0}^{N}g_{2j}T_{2j}(x),\qquad N=9^m,
\]
the finite residue is \(F_N(g)=\Pi_N(\mathcal Rg-g)\). If \(A_N\) approximates
\((DF_N(\bar g))^{-1}\), the quantities that must be certified are
\begin{align*}
Y&=\|A_NF_N(\bar g)\|,\\
Z_0&=\|I-A_NDF_N(\bar g)\|,\\
Z_1(r)&=\sup_{\|h\|\le r}
\|A_N(DF_N(\bar g+h)-DF_N(\bar g))\|.
\end{align*}
An inequality
\(Y+(Z_0+Z_1(r))r<r\) encloses a unique fixed point. The later spectral block
must isolate an unstable direction and contract the complement. The
use of \(9^m\) thus enters as a refinement law of the certificate, not as
decoration numerological.

\begin{corollary}[Exact convergence of mode thirty]
Mode \(30\) admits three simultaneous and typologically
distinct realizations:
\begin{align*}
30&\longmapsto(3\cdot30,4\cdot30)=(90,120)
&&\text{in the constitutive completion},\\
30&\longmapsto30^2-1=899=29\cdot31
&&\text{in the critical normalization},\\
30&\longmapsto
M_{30}=\begin{pmatrix}30&899\\1&30\end{pmatrix}
&&\text{in the hyperbolic Pell transformation}.
\end{align*}
The three identities are exact and are obtained without prescribing target values. The conserved invariant is the central integer \(30\) together with its orientation: extension \(3\to4\), quadratic defect \(-1\), and a unit of norm one. What is not asserted is the equality between the vacuum operator and the duplication renormalizer.
\end{corollary}

\begin{remark}[Negative control]
A second eigenvalue outside the unit disk, a ball of radii without solution, loss of uniqueness, or tangency of the HMT family to the stable variety refutes the universal promotion. None refutes the exact germ or the Schwarzian theorem.
\end{remark}

\begin{proposition}[Main Result]
The structural significance of the doubling constant lies in the
cancellation in dodecaphase form of \(\pi\), in the identity
\(899=30^2-1\), in the closed trigonometric form, and in its proved membership
in the analytic class on which the renormalizer acts. Its
decimal expansion constitutes exclusively the terminal evaluation.
\end{proposition}

\subsection*{Logical Scope and Domain of Validity}
Profile, Taylor expansion, normalizer, Pell unit, unimodality, and Schwarzian: \emph{Exact internal theorem}. Approximants through level eight: \emph{Persistent finite certificate}. Levels 11--14: \emph{New numerical certificate pending durability}. Hyperbolicity and transversality of the renormalizer: \emph{Open program with executable falsification criterion}.

